% ============================================================================
\clearpage
\section{Notación}\label{ap:notacion}
% ============================================================================

\begin{table}[H]
\centering\small
\caption{Símbolos principales. Algunos se reutilizan con significado local ($\alpha$, $\varepsilon$, $r$);
el contexto de cada sección los desambigua.}
\label{tab:notacion}
\begin{tabularx}{\linewidth}{@{}l Y@{}}
\toprule
\encab{Símbolo} & \encab{Significado}\\
\midrule
$S_t$, $F_{t,\tau}$, $D(t,T)$ & Precio del subyacente, forward al plazo $\tau$ y factor de descuento\\
$K$, $T$, $\tau$ & Strike, vencimiento y plazo constante de la distribución implícita (30 días)\\
$k=\ln(K/F)$, $w=\sigma^2T$ & Log-moneyness y varianza implícita total\\
$\theta_T$, $\rho$, $\varphi$, $\eta$, $\gamma$ & Parámetros SSVI: varianza ATM, asimetría, curvatura y función de potencia\\
$b_i$, $a_i$, $m_i$, $s_i$ & Bid, ask, mid y medio spread de la cotización $i$\\
$q^{\Q}(K)$, $g(k)$ & Densidad implícita del precio y función de Durrleman\\
$x_t(u)$, $u$ & Cuantil de $\ln(S_{t+\tau}/F_{t,\tau})$ bajo $\Q_t$ y nivel de probabilidad\\
$W_1$, $W_2$, $\Delta x_t(u)$ & Distancias de Wasserstein y cambio firmado de cuantiles\\
$\phi_j(u)$, $s_{tj}$ & Autofunciones y puntajes de la PCA funcional\\
$f_t$, $A$, $Q$, $R_t$ & Estado, transición, ruido de estado y ruido de observación del filtro\\
$\mathcal I^{base}_t$, $\epsilon_t(u)$ & Información base e innovación de la distribución implícita\\
$Y_{t,h}$, $d_{t,h}$ & Objetivo a horizonte $h$ y diferencial de pérdida entre modelos\\
$\rho_\tau$, CRPS, Brier & Pérdida cuantílica y reglas de puntuación\\
$r_t$, $H$ & Longitud de corrida y probabilidad de riesgo de BOCPD\\
$\alpha_t$, $\gamma$ & Nivel efectivo y paso del conformal adaptativo\\
$w$, $\kappa$, $\lambda$, $\mathcal C_t$ & Pesos, costos proporcionales, aversión al riesgo y conjunto factible\\
$\CVaR_\alpha$, $\mathcal R_\delta$ & Valor en riesgo condicional al nivel $\alpha$ y región casi óptima\\
$\mathcal B_\varepsilon(\widehat P_t)$ & Bola de Wasserstein de radio $\varepsilon$ alrededor de la distribución estimada\\
\bottomrule
\end{tabularx}
\end{table}

% ============================================================================
\section{Glosario}\label{ap:glosario}
% ============================================================================

\begin{description}[leftmargin=0pt,itemsep=3pt,font=\sffamily\bfseries\small]
\item[Arbitraje estático.] Combinación de opciones de un mismo instante que asegura ganancia sin riesgo. Sus
dos formas son de \emph{calendario} (entre vencimientos) y de \emph{mariposa} (entre strikes).
\item[Breeden--Litzenberger.] Identidad que obtiene la densidad implícita como segunda derivada del precio
de las calls respecto del strike.
\item[BOCPD.] Detección bayesiana en línea de puntos de cambio mediante la posterior de la longitud de corrida.
\item[Conformal adaptativo.] Método de intervalos que ajusta su nivel efectivo según los errores recientes y
garantiza la tasa de error promedio a largo plazo.
\item[CRPS.] Regla de puntuación de pronósticos de distribución; integra la pérdida cuantílica en todos los niveles.
\item[CVaR.] Promedio de las pérdidas en la cola especificada (por ejemplo, el 5\% peor) bajo los escenarios supuestos.
\item[Desamericanización.] Conversión de precios de opciones americanas en equivalentes europeos mediante un
modelo con ejercicio anticipado y dividendos.
\item[Embargo y purga.] Eliminación de observaciones cuya etiqueta cruza una frontera temporal entre tramos, más
un margen de seguridad.
\item[Idempotencia.] Propiedad de una orden cuyo reintento con la misma clave no crea duplicados.
\item[Innovación.] Parte del cambio de la distribución implícita no explicada por la información base.
\item[PIT.] Valor $F(y)$ del pronóstico en el resultado observado; uniforme si el pronóstico está calibrado.
\item[$\Q$ y $\P$.] Distribución implícita en los precios de las opciones y distribución física de los
rendimientos.
\item[Región casi óptima.] Conjunto de carteras factibles cuyo valor estimado está dentro de la incertidumbre
del óptimo; base de la zona de no operación.
\item[Sharpe deflactado.] Probabilidad de que el Sharpe verdadero supere el máximo esperado por azar dado el
número de variantes probadas.
\item[SSVI.] Parametrización de la superficie de varianza implícita con condiciones suficientes sencillas de
ausencia de arbitraje estático.
\item[Transporte óptimo.] Forma de mover la masa de una distribución a otra con costo mínimo; en una dimensión
es el reordenamiento monótono de cuantiles.
\end{description}

% ============================================================================
\section{Decisiones pendientes}\label{ap:pendientes}
% ============================================================================

\begin{table}[H]
\centering\small
\caption{Decisiones que parametrizan el sistema y dónde se usan.}
\label{tab:pendientes}
\begin{tabularx}{\linewidth}{@{}c l Y@{}}
\toprule
& \encab{Decisión} & \encab{Dónde se usa}\\
\midrule
\iconopendiente & Universo definitivo & Datos, escenarios conjuntos, límites de concentración\\
\iconopendiente & Moneda de los 50,000 & Costos, conversión, benchmark comparable\\
\iconopendiente & Posiciones actuales & Comparación mantener frente a modificar; zona de no operación\\
\iconopendiente & Horizonte de decisión principal & Objetivos, pérdidas, etiquetas y purga\\
\iconopendiente & Presupuesto y proveedor de datos & Historia disponible, latencia y horario de decisión\\
\iconopendiente & Modelo de costos & Programa lineal y zona de no operación\\
\iconopendiente & Presupuesto de riesgo y límites & Conjunto factible $\mathcal C_t$ y reducciones obligatorias\\
\iconopendiente & ¿Se gestionarán opciones? & Rama de revaluación completa\\
\iconopendiente & Umbrales de la regla de avance & Deben fijarse antes de abrir la prueba final\\
\bottomrule
\end{tabularx}
\end{table}

% ============================================================================
\section{Reproducibilidad de las figuras}\label{ap:reproducibilidad}
% ============================================================================

Todas las gráficas se generan con datos sintéticos, semillas fijas y el paquete de referencia
\texttt{quantileflow/} del repositorio. Se dibujan con matplotlib en su estilo por defecto, una figura por
gráfica y con textos en inglés; los pies, en español, explican cada una. Desde la raíz del repositorio:

\begin{lstlisting}[style=registro]
pip install -r requirements.txt
python -m pytest                              # pruebas de propiedades del núcleo
python docs/propuesta/figuras_src/generar.py  # regenera las 62 figuras PNG
make pdf                                      # compila este documento
\end{lstlisting}

\begin{table}[H]
\centering\small
\caption{Origen de las figuras.}
\label{tab:origen}
\begin{tabularx}{\linewidth}{@{}l l Y@{}}
\toprule
\encab{Figuras} & \encab{Script} & \encab{Mundo sintético}\\
\midrule
\texttt{s1\_*}, \texttt{s2\_*} & \texttt{fig\_s1\_s2.py} & Calendario de 30 días; cadena SSVI con spreads y bid nulo\\
\texttt{s3\_*} & \texttt{fig\_s3.py} & Panel de 260 sesiones: mercado y acción A, choque común y evento idiosincrático\\
\texttt{s4\_*}, \texttt{s5\_*} & \texttt{fig\_s4\_s5.py} & Supuesto $\Q$/$\P$; regresión cuantílica; cambio de varianza; desplazamiento de volatilidad\\
\texttt{s6\_*}, \texttt{s7\_*}, \texttt{s9\_*} & \texttt{fig\_s6\_s7.py} & Cópulas; cartera de cuatro activos; remuestreo de óptimos; nulas de Sharpe\\
\bottomrule
\end{tabularx}
\end{table}

Las pruebas verifican, entre otras propiedades, que la superficie de referencia cumple las condiciones
suficientes de SSVI, que el contraejemplo SVI tiene $g<0$, que la densidad por Breeden--Litzenberger
coincide con la analítica e integra a uno con media igual al forward, que $W_2$ de una traslación es la
traslación, que el CVaR del programa lineal coincide con el empírico, que costos altos producen no
operación, que el conformal adaptativo cumple su cota de largo plazo y que el máximo Sharpe esperado
coincide con la simulación.
